\documentclass[11pt,a4paper]{article}
\usepackage[T1]{fontenc}
\usepackage[utf8]{inputenc}
\usepackage{lmodern}
\usepackage[margin=25mm,headheight=14pt]{geometry}
\usepackage{amsmath,amssymb,amsthm,mathtools}
\usepackage{booktabs,longtable,array,enumitem,microtype}
\usepackage[dvipsnames]{xcolor}
\usepackage{xurl}
\usepackage{hyperref}
\usepackage{fancyhdr}
\definecolor{inkblue}{HTML}{183E5A}
\hypersetup{colorlinks=true,linkcolor=inkblue,citecolor=inkblue,urlcolor=inkblue,
 pdftitle={Three explicit Diophantine forms of ZFC consistency},
 pdfauthor={Research synthesis for ProveIt},pdfsubject={Nine, twelve and twenty-seven variable consistency equations}}
\pagestyle{fancy}
\fancyhf{}
\fancyhead[L]{\small\textsc{Diophantine forms of ZFC consistency}}
\fancyhead[R]{\small 22 October 2026}
\fancyfoot[C]{\thepage}
\renewcommand{\headrulewidth}{0.3pt}
\setlength{\parskip}{3pt}
\setlength{\emergencystretch}{2em}
\setlist{itemsep=3pt,topsep=5pt}
\allowdisplaybreaks[3]
\newtheorem{theorem}{Theorem}[section]
\newtheorem{lemma}[theorem]{Lemma}
\newtheorem{proposition}[theorem]{Proposition}
\theoremstyle{remark}
\newtheorem{remark}[theorem]{Remark}
\newcommand{\NN}{\mathbb N}
\newcommand{\ZZ}{\mathbb Z}
\newcommand{\Con}{\operatorname{Con}}
\newcommand{\ZFC}{\mathrm{ZFC}}
\newcommand{\pin}{\texttt{f0ae4e0991ca44a7bb4e5671d42e33d775e48035}}
\newcommand{\repofile}[2]{\href{https://github.com/VladimirReshetnikov/ProveIt/blob/f0ae4e0991ca44a7bb4e5671d42e33d775e48035/#1}{#2}}
\title{\textbf{Three explicit Diophantine forms\\of ZFC consistency}\\[6pt]
 \large Nine nonnegative unknowns, twelve positive unknowns,\\
 and twenty-seven unrestricted integer unknowns}
\author{Research synthesis prepared for the ProveIt repository}
\date{22 October 2026\\[3pt]\small Repository comparison frozen at \texttt{f0ae4e0991ca}}

\begin{document}
\maketitle
\begin{abstract}
We combine the exact blank-input K215 consistency machine, the native
tableau/Pell compiler, the four-cell T4 coefficient recipe, the existing
32-variable quartic, and the constructive nine-unknown reduction in the
repository's formalization of Jones's 1982 work.  This yields a fully
specified ordinary polynomial in nine nonnegative integers whose lack of
solutions is provably equivalent in ZFC to $\Con(\ZFC)$.  Every coefficient
is fixed by a bounded recipe; no recursively enumerable index is left
unspecified.  Smaller unconditional radical weights reduce the degree of
this particular specialization to $161844\cdot5^{33}+484$.

Two further constructions improve the direct arithmetic family.  A new
four-root combination gives twelve positive unknowns, degree $1609$ and
$151$ binary arithmetic operations.  A signed three-root combination gives
twenty-seven unrestricted integer unknowns, degree $1004$ and $137$
operations, or four equations of maximum degree $952$ and $131$
operations.  These are different size tradeoffs.  They do not improve the
existing $67$-operation record and are not asserted globally minimal.

The new arguments are ordinary mathematical proofs, with exact polynomial
data, reproducible finite audits and an independent mathematical review.
They import the existing compiler and K215 calibration proofs.  The archive
does not claim a new end-to-end Lean/Rocq build or an internally encoded
ZFC derivation.
\end{abstract}

\clearpage
\begingroup
\small
\setlength{\parskip}{0pt}
\tableofcontents
\endgroup
\clearpage

\section{Results, conventions and comparison}
\label{sec:results}

Throughout, $\NN=\{0,1,2,\ldots\}$ and
$\ZZ_{>0}=\{1,2,3,\ldots\}$.  A displayed intermediate expression is a
polynomial abbreviation unless it is explicitly declared an unknown.
All exponents in the final equations are fixed nonnegative integers.
Large finite sums, fixed coefficient tables and fixed powers are
descriptions of ordinary polynomials, not extra existential quantifiers.

\begin{theorem}[Main results]
\label{thm:main}
The formulas and fixed coefficient recipes in this report define
polynomials $\mathcal F_9$, $\mathcal P_{12}$ and $\mathcal P_{27}$ with
integer coefficients such that ZFC proves each of
\begin{align*}
 (\nexists \mathbf x\in\NN^9:\mathcal F_9(\mathbf x)=0)
      &\ \longleftrightarrow\ \Con(\ZFC),\\
 (\nexists \mathbf x\in\ZZ_{>0}^{12}:\mathcal P_{12}(\mathbf x)=0)
      &\ \longleftrightarrow\ \Con(\ZFC),\\
 (\nexists \mathbf x\in\ZZ^{27}:\mathcal P_{27}(\mathbf x)=0)
      &\ \longleftrightarrow\ \Con(\ZFC).
\end{align*}
The alternative systems listed below have the same existence predicates.
The claim is a provable equivalence, with no assumption of consistency
needed for its proof.
\end{theorem}

Here $\Con(\ZFC)$ means that there is no finite syntactic ZFC derivation
of falsity, with the conventional effective first-order presentation
used by the K215 calibration.  Equivalent standard proof encodings give
the same assertion by their effective translations.

\begin{center}
\small
\begin{tabular}{@{}llrrrr@{}}
\toprule
Construction & Domain & Unknowns & Equations & Exact degree & Operations\\
\midrule
Nine, new weights & $\NN$ & 9 & 1 & $161844\cdot5^{33}+484$ & 20,531\\
Direct twelve & $\ZZ_{>0}$ & 12 & 1 & 1609 & 151\\
Direct twelve & $\ZZ_{>0}$ & 12 & 3 & 1600 & 147\\
Signed, lower weight & $\ZZ$ & 27 & 1 & 1004 & 137\\
Signed, lower weight & $\ZZ$ & 27 & 4 & 952 & 131\\
Signed, full weight & $\ZZ$ & 27 & 1 & 1116 & 136\\
Signed, full weight & $\ZZ$ & 27 & 4 & 1064 & 130\\
\bottomrule
\end{tabular}
\end{center}

For a system, degree means the maximum total degree of a residual.  An
operation is one binary addition, subtraction or multiplication in the
supplied shared arithmetic circuit, including squaring and multiplication
by numerals.  Equality comparisons, register reuse and supplied fixed
coefficient leaves are free, as in the repository's comparisons.  The
nine-variable circuit expands all fixed powers into multiplications.
The counts describe these emitted circuits, not a lower bound on every
possible evaluator.

\subsection{What becomes simpler}

The baseline is the repository at \pin.  Its wave-2 comparison
\cite{collection} reports thirteen positive unknowns for the smallest
displayed direct construction, with degree $571$ and $84$ operations
(or degree $515$ and $85$ operations).  Its best listed unrestricted
integer count is twenty-nine \cite{integer29}.  The Hilbert-tenth README
at this pin explicitly distinguishes the formal existence theorem for
nine unknowns from the then-open concrete consistency specialization.
Our nine-variable construction fills that fixed-instance specification
gap; the general nine-unknown theorem itself is classical \cite{jones1982}
and is not new here.

The direct twelve-variable construction deletes another positive
coordinate by combining the main Pell root with the three roots already
amalgamated in package~09.  The twenty-seven-variable construction proves
that the three-root amalgamation survives the signed domain of
package~14, including two exceptional zero cases that the positive
argument could not discard automatically.

The earlier points still matter.  A degree-$121$, fourteen-positive-variable
polynomial costs $76$ operations; a fifteen-variable polynomial costs
$67$ operations at degree $155$; the $32$-variable quartic is a sum of
squares of $23$ quadratics.  There is also a fifteen-variable six-equation
system of degree $22$ at $61$ operations, and a degree-$28$ system at
$60$ operations \cite{collection}.  Thus ``simplest'' has no single
answer without a choice of variable count, degree, circuit size,
coefficient height, or ease of proof.  This report adds specific points
to that comparison rather than claiming one dominates all the others.

\subsection{Coefficient size and the domain matter}

The basic T4 radix is already
\[
 B=2^{5^{18}}=2^{3\,814\,697\,265\,625}.
\]
The binary expansion of $\beta=B-1$ alone has
$3\,814\,697\,265\,625$ bits, approximately $477$ billion bytes.
The Jones index constants are much larger.  We specify these integers
exactly by bounded formulas and supply modular evaluators; we do not
pretend that their decimal expansion would be a useful artifact.

The domain of $\mathcal F_9$ is essential: its last nonnegative coordinate
enforces a strict inequality.  Using the same polynomial over all of
$\ZZ^9$ is not justified.  Nine strictly positive coordinates are possible
by substituting $x_j-1$ for each nonnegative coordinate, without changing
degree.  Converting all nine coordinates to unrestricted integers with
three-square representations gives another $27$-integer form, but retains
the huge degree.  The separate signed construction below is far smaller
under the degree and circuit measures.

\section{The logical and compiler interface}
\label{sec:interface}

\subsection{The exact machine and consistency predicate}

Let $K215$ be the repository's dense two-symbol, $215$-working-state
machine, started in state $0$ on the blank tape, with a separate halt
state $215$.  The SHA-256 of its literal dense table is
\begin{center}\small\ttfamily
fdbcef0825886035c3091c8537a69535bc96f9e265f0ae45863543d40ae70090.
\end{center}
The source theorem
\texttt{KEndpoint215.blank\_halts\_iff\_zfc\_refutation} has endpoint
\[
 \operatorname{BlankHalts}(K215)\quad\longleftrightarrow\quad
 \operatorname{BProv}(\ZFC\mathrm{ax}_s,[],\bot).
\]
The right side is finite syntactic refutability.  The ordinary calibration
proof explicitly explains how the machine simulation, initialization,
source-program equivalence and dominated-calculus calibration compose
to a proof formalizable in ZFC \cite{k215ordinary,k215formal}.
In the formal source the shared calibration passes through ZF
refutability and the constructible-choice interpretation relating ZF
and ZFC refutations.  There is no consistency premise in the endpoint.

The published K215 records report successful Lean and Rocq endpoint
builds for this table.  They also explicitly distinguish these external
proof-assistant results from a derivation encoded in ZFC's own syntax.
Our use of ``ZFC proves'' rests on the ordinary calibration argument
formalizable in ZFC, together with the arithmetic proofs here.  It is
not an inference that a Lean theorem, merely by existing, supplies a
ZFC proof object.  These builds were not rerun for this report.

\subsection{One complete coefficient tuple}

All constructions here use the complete T4 tuple of package~11
\cite{t4}.  Write its four fixed polynomial coefficients as
\[
 \beta=B-1,\qquad K_0=D_C+BD_R,\qquad M_C,\qquad
 M_F=M_F^{\rm native}+B-1.
\]
Section~\ref{sec:t4-recipe} defines them without an unbounded search.
In particular, the fourth coefficient is the \emph{shifted} mask.
The names \texttt{Bm1} and \texttt{Kconstant} in inherited circuits mean
$\beta$ and $K_0$, respectively.  They do not denote additional inputs.
The root \texttt{COEFFICIENT\_BINDINGS.json} makes these aliases explicit.

The native-window compiler converts a marked cyclic T4 word into
arithmetic constraints.  The overlap relations pull back the cyclic
word to a globally valid plane even for an arbitrary stride: no
coprimality or divisibility assumption on that stride is needed.
The marked initialization tile then identifies a finite blank-input
computation strip.  Conversely, every finite halt has sufficiently
wide blank margins and an arbitrarily padded halted tail; independent
width and height powers of five give the required cyclic presentation.
The mask, population and Pell lemmas turn these word constraints into
the fixed four-port arithmetic interface \cite{t4,pellinterface}.

Package~11 proves that this entire T4 tuple satisfies that interface.
The low-degree lift and the root eliminations use its interface
identities and inequalities, not a particular historical ordering of
larger tiles.  We therefore substitute the tuple as a whole.  Mixing
one port from T4 with a historical mask or tile ordering would not be
covered by the proof.

\subsection{What is imported and what is proved here}

The imported mathematical components are: the K215 calibration; the
native/Pell compiler including its positive completion; the T4 adapter;
the ordinary strong-rank and signed-core arguments; and the constructive
quartic-to-nine reduction.  Selected exact source files and a source
manifest accompany the report.  Their transitive historical proof
closure is not claimed to be a newly self-contained formal project.

The new work is the explicit normalized quartic and bounded index
specialization; smaller unconditional weights in the nine-variable
finalizer; the twelve-variable four-root proof; and the signed
twenty-seven-variable proof.  Algebraic generators, exact identity
checks, degree certificates and modular evaluations support these
constructions.  A finite check cannot replace a compiler theorem about
arbitrary runs or the universal arithmetic inequalities used below.

For any of the constructions the proof chain is
\[
 \text{a solution in the stated domain}
 \ \longleftrightarrow\
 \text{a finite blank halt of }K215
 \ \longleftrightarrow\
 \text{a finite ZFC refutation}.
\]
Negating this chain is exactly the required consistency equivalence.

\section{A concrete normalized quartic and its bounded index}
\label{sec:quartic}

\subsection{The twenty-three quadratic equations}

We use package~09's ordinary-strong $32$-variable lift
\cite{quartic32}.  Its unknowns, in the exact coefficient-file order, are
\begin{align*}
\mathbf v=(
&J,F,\alpha,z,f,h,i,o,j,s,w,\tau,\eta,\zeta,y,\gamma,\\
&N,E,P,c,a,X,D,C_2,G,L,F_2,V,g,I,B_{\rm aux},\Delta)
\ \in\ZZ_{>0}^{32}.
\end{align*}
The affine abbreviations are
\[
 q=\beta J+1,\qquad k=\eta+\zeta,\qquad H=4a+3,
 \qquad C_w=q-F-\alpha.
\]
The residuals $r_j$, all of total degree at most two in these unknowns,
are the following.  The four fixed ports have degree zero.

\begin{longtable}{@{}r@{\quad}l@{}}
\toprule
$j$ & $r_j(\mathbf v)$\\\midrule
\endfirsthead
\toprule $j$ & $r_j(\mathbf v)$\\\midrule
\endhead
1 & $L-q^2+1$\\
2 & $N-q(L+1)$\\
3 & $X-wq$\\
4 & $a-E-sN$\\
5 & $E-X(a-E)$\\
6 & $c-k(a-E)-\eta$\\
7 & $P-E(c-\eta)-k$\\
8 & $\tau^2-P(P-k)-1$\\
9 & $D-X-ac-\gamma H$\\
10 & $\Delta-a^2-H$\\
11 & $C_2-c^2$\\
12 & $D^2-\Delta C_2-1$\\
13 & $k-jc+V-1-hE$\\
14 & $G-(q-1)(q-F)-\alpha$\\
15 & $jc-V-LG-(M_C+qM_F)J$\\
16 & $(K_0+w)C_w+q-F-z(q-1)-1$\\
17 & $F_2-f^2$\\
18 & $V-of+c$\\
19 & $g-y^2+V^2$\\
20 & $I-iC_2$\\
21 & $B_{\rm aux}-I^2$\\
22 & $B_{\rm aux}-\Delta(F_2-1)$\\
23 & $y^2-B_{\rm aux}g-1$\\
\bottomrule
\end{longtable}

This is already a useful low-degree answer to the original question:
the simultaneous equations $r_1=\cdots=r_{23}=0$, or the single quartic
$\sum r_j^2=0$, have a positive solution exactly when K215 halts.
It is an inherited construction, included to make the nine-variable
specialization fully determined.

For orientation, the first six rows recover
$N=q^3$, $Y=a-E=sq^3$, $X=wq$, $E=XY$, and $c=kY+\eta$.
Rows 7--8 give the factored first Pell norm.  Rows 9--12 give the main
Pell norm.  Rows 13--16 encode the packed index and transport equation.
The remaining rows give the ordinary strong relation
$(ic^2)^2=\Delta(f^2-1)$ and the auxiliary norm.
This ordinary strong relation is deliberately distinguished from
$f^2-\Delta i^2c^4=1$, used by the normalized direct constructions.
The ordinary strong-rank lemma supplies soundness; positive canonical
witnesses supply every lifted register, including the positive gap
$g=y^2-V^2$.  These projection and extension claims are proved in the
bundled low-degree proof, not inferred merely from the number of rows.

\subsection{Normalization is explicit}

Introduce $32$ nonnegative coordinates $n_1,\ldots,n_{32}$ and a dummy
input $x$ which does not occur in the polynomial.  Define
\begin{equation}
 P_{32}(x,\mathbf n)=\sum_{j=1}^{23}r_j(\mathbf n+\mathbf 1)^2.
 \label{eq:quartic}
\end{equation}
Because a sum of integer squares is zero exactly when each summand
vanishes, its nonnegative zeros are precisely the positive solutions
of the quadratic system.  It has degree exactly four: the coefficient
of $n_\tau^4$ is one.

At the all-ones positive tuple, the fifth residual is
$1-1(1-1)=1$.  Consequently
\[
 P_{32}(x,0,\ldots,0)\ge1
 \quad\text{for every }x\in\NN.
\]
This discharges the normalization hypothesis of the Jones construction
without an extra variable or a case split.  The exact upstream sparse
file has SHA-256
\begin{center}\small\ttfamily
149eacd479896a113d869c5b74163efa4dab2cabaca8b581a6405cc9f580ca60.
\end{center}

\subsection{Every index constant is specified}

Write $\nu=32$, $L=5^{33}$, and let
$\alpha=(\alpha_0,\ldots,\alpha_{32})\in\NN^{33}$ with
$|\alpha|\le4$.  Define
\begin{align*}
 p_\alpha&=[x^{\alpha_0}n_1^{\alpha_1}\cdots n_{32}^{\alpha_{32}}]P_{32},\\
 c_\alpha&=(4-|\alpha|)!\prod_{j=0}^{32}\alpha_j!\;p_\alpha,
 &e(\alpha)&=\sum_{j=0}^{32}\alpha_j5^j.
\end{align*}
The slack factorial is essential.  Merely using the ordinary monomial
coefficient $p_\alpha$ would not satisfy the source coding identity.
There are $\binom{37}{4}=66\,045$ eligible multi-indices.  The shifted
quartic has $346$ nonzero witness-coefficient polynomials, comprising
$508$ terms when the four fixed ports are retained symbolically.
All nonzero records have $\alpha_0=0$, but the baseline sum below still
includes the multi-indices with a nonzero input exponent.

Here is a deterministic bound with no coefficient search.  For each raw
residual, first expand the affine abbreviations in the table and write
\[
 r_j(\mathbf v)=\sum_{a,b}d_{j,a,b}\,
       \mathbf v^a(\beta,K_0,M_C,M_F)^b,
 \qquad d_{j,a,b}\in\ZZ.
\]
Set
\[
 \rho_j=\sum_{a,b}|d_{j,a,b}|\,2^{|a|}
                  (\beta,K_0,M_C,M_F)^b,\qquad
 s_\star=2+24\sum_{j=1}^{23}\rho_j^2,
 \qquad z_\star=2^{s_\star}.
\]
All four ports are positive integers.  The factor $2^{|a|}$ bounds
the coefficient $\ell^1$ norm after the shift $\mathbf v=\mathbf n+1$;
the norm is submultiplicative.  Each factorial multiplier is at most
$4!=24$.  Therefore
\[
 \sum_{|\alpha|\le4}|c_\alpha|
       \le24\sum_j\rho_j^2=s_\star-2,
 \qquad |c_\alpha|+1<2^{s_\star}=z_\star.
\]
The explicitly expanded bound is only a twenty-term polynomial:
\begin{align*}
 s_\star={}&79250+960M_F+96M_F^2+960M_C+192M_CM_F+96M_C^2\\
 &+2880K_0+600K_0^2+9504\beta+1920\beta M_F
       +384\beta M_F^2+384\beta M_CM_F\\
 &+3552\beta K_0+480\beta K_0^2+6048\beta^2
       +384\beta^2M_F^2+960\beta^2K_0+96\beta^2K_0^2\\
 &+1920\beta^3+768\beta^4.
\end{align*}
Finally define the fixed digit polynomial and constants
\begin{align*}
 e_\star(T)&=\sum_{|\alpha|\le4}(z_\star+c_\alpha)T^{L-e(\alpha)},\\
 u_\star&=\sum_{j=1}^{32}(2z_\star)^{5^j},
 &y_\star&=e_\star(2z_\star).
\end{align*}
All exponents are nonnegative since
$e(\alpha)\le4\cdot5^{32}<L$.
The coefficient bound makes every digit positive; in particular
$z_\star,u_\star,y_\star$ are natural numbers.  These definitions satisfy
the source predicate \texttt{Index 32 $P_{32}$ $z_\star$ $u_\star$ $y_\star$}
exactly \cite{nineformal}.  The value $u_\star$ is part of the proof
interface but is not a free leaf of the final polynomial.

\input{nine/nine_section.tex}
\input{direct12/section.tex}
\input{integer27/integer27_section.tex}

\section{The fixed T4 coefficient recipe}
\label{sec:t4-recipe}

This section freezes one complete instance of the compiler.  It
reproduces the bounded definition proved by package~11 \cite{t4}, with
separate notation for the native and shifted masks.  The full ordered
tile list, literal K215 transitions, predicate and exact coefficient
metadata are included in the archive.  Neither the tile enumeration
nor the coefficient definition asks whether K215 halts.

\subsection{The finite local language}

There are $439$ cell symbols.  Headless tape bits have codes $0,1$;
a bit $b\in\{0,1\}$ with head state $q\in\{0,\ldots,215\}$ has
code $2(q+1)+b$.  The special cells are
\[
 L=434,\quad Q=435,\quad R=436,\quad V=437,\quad H=438.
\]
Here $L,Q,R$ are initialized phases with payloads blank/headless,
blank/state-$0$, blank/headless.  $V$ is a vertical boundary and $H$ a
horizontal boundary.  They carry no tape payload.  Only the completed
halt transition stays in place: $(215,b)\mapsto(b,0,215)$.  The actual
first K215 action is $(0,0)\mapsto(1,+1,93)$; the transition entering
halt retains its actual write and movement.

A tile is a quadruple $(l,c,r,n)$ of left, center, right and next-time
center cells.  The exact predicate is the supplied
\texttt{upstream/t4/bridge/check\_t4.py}.  In mathematical terms:
\begin{enumerate}
\item At most one of $l,c,r$ has a head; the initialized phase $Q$
counts as a head.
\item $c$ is $V$ exactly when $n$ is $V$.  Adjacent $c,r$ cannot both
be $V$, and $c,n$ cannot both be $H$.
\item If neither $c$ nor $r$ is $V$, they agree on whether they are $H$.
For a center $V$, there is no further update requirement.
\item A center $H$ has next-time center in $\{L,Q,R\}$.
For an initialized interior center, the left boundary requires $L$,
the right boundary requires $R$, and away from a right boundary the
phase edge is one of $LL,LQ,QR,RR$.
\item An interior head may not move into an adjacent $V$.  If the next
center is $H$, the current center is headless or has halt state $215$.
\item Otherwise the next center is ordinary, with no initialized phase,
and is the exact three-neighbor deterministic successor.  A neighbor $V$
is blank/headless; a neighbor $H$ cannot supply this update.  The new
bit is the center's written or stable bit, and incoming heads are given
by the completed transition table.  Multiple arrivals are forbidden.
\end{enumerate}

Place $S=(434,435,436,1)$ first and $E=(438,0,0,438)$ second, then
list the other passing tuples lexicographically.  Both distinguished
tiles pass the local predicate.  The second cannot occur in a valid
plane: its left $H$ would have a non-boundary, non-$H$ right neighbor,
contradicting horizontal boundary propagation.  Forbidding this tile
therefore excludes no genuine computation.

There are exactly $37\,075$ tiles.  The independent count by center
class is
\begin{center}
\begin{tabular}{@{}lr@{\qquad}lr@{}}
\toprule
Center class & Tiles & Center class & Tiles\\\midrule
Headless ordinary 0 & 9,148 & Headless ordinary 1 & 9,148\\
Working-head symbols & 8,600 & Halt-head symbols & 110\\
$L$ & 888 & $Q$ & 4\\
$R$ & 1,750 & $V$ & 4,793\\
$H$ & 2,634 & Total & 37,075\\\bottomrule
\end{tabular}
\end{center}
The ordered list has SHA-256
\begin{center}\small\ttfamily
9c2e473b5dbc5b3dc1468a31aefdbec6a901d6ba1e6b1ec9212698d41e1d84b1.
\end{center}

For a cyclic sequence of tiles $w_i$ with stride $h$, the three overlaps
are $w_i(r)=w_{i-1}(c)$, $w_i(c)=w_{i-1}(l)$ and
$w_i(n)=w_{i-h}(c)$.  With $a_i=w_i(c)$ this gives
$w_i=(a_{i+1},a_i,a_{i-1},a_{i-h})$, hence a valid pullback plane
$a_{-x-ht}$.  This is why residue collisions and arbitrary strides do
not invalidate soundness.  The marked initialized row lies between
vertical boundaries and the following horizontal boundary requires a
halted last interior head.  For completeness, sufficiently wide blank
margins make the added one-head filter harmless, including helical
seams.  Independent width and height powers of five remain available.

\subsection{Literal coefficient formulas}

Put
\[
 A=2^{16},\quad\rho=2^{390625},\quad
 B=\rho^{9765625}=2^{5^{18}},\qquad
 \mu=A-2+\sum_{j=1}^{18743}A^j.
\]
For $0\le u\le438$, use copy positions
\[
 e_l(u)=40587+u,\quad e_c(u)=40148+u,\quad
 e_r(u)=39709+u,\quad e_n(u)=38831+u.
\]
Let $\mathcal E$ contain selector positions $0,\ldots,37074$, these
$1756$ copy positions, the dummy $42343$, and the four anchors
$43661,130983,392949,1178847$.  Its size is $38\,836$.
For the ordered tile $w_i=(l_i,c_i,r_i,n_i)$ define
\begin{align*}
 v_i&=1+A^{1+l_i}+A^{440+c_i}+A^{879+r_i}+A^{1318+n_i}
               +\sum_{t=1}^4 A^{1756+t},\\
 v_{e_l(u)}&=A^{1+u},&v_{e_c(u)}&=A^{440+u},\\
 v_{e_r(u)}&=A^{879+u},&v_{e_n(u)}&=A^{1318+u},\\
 v_{3^{t-1}43661}&=A^{1756+t}\quad(1\le t\le4),&v_{42343}&=0.
\end{align*}
Then
\begin{align*}
 D_C={}&\rho^{1317}+\rho^{2228468}+\rho^{349288}+\rho^{1047864}\\
 &+\sum_{e\in\mathcal E}v_e
       (\rho^{4586163-e}+\rho^{6943861-e}),\\
 D_R={}&\rho^{2228029},\\
 M_F^0={}&\mu(\rho^{4586163}+\rho^{6943861})
       +\sum_{u=0}^{438}\rho^{40148+u}\\
 &+\sum_{u=0}^{438}
       (\rho^{2228029+40587+u}+\rho^{2228029+40148+u})
       +\rho^{392949}+\rho^{1178847},\\
 M_C={}&B-1-\sum_{\substack{e\in\mathcal E\\e\ne1}}\rho^e-2\rho^{42343},\\
 M_F^{\rm native}={}&M_F^0+4.
\end{align*}
The fixed ports used in every final equation are
\begin{equation}
 \boxed{\beta=B-1,\quad K_0=D_C+BD_R,\quad M_C,
       \quad M_F=M_F^0+4+B-1.}
 \label{eq:fixed-ports}
\end{equation}
These are positive integers satisfying the compiler's mask and size
conditions.  Package~11 proves the support geometry, carry bounds,
population census and five-adic index adjustment.  Our replayed finite
checks verify their finite data; the arbitrary-length coding proof
remains the imported mathematical theorem.

\section{Verification, reproducibility and scope}
\label{sec:verification}

\subsection{Executed checks}

The supplied Python programs use the standard library.  The principal
checks performed for this report are:
\begin{itemize}
\item Exact sparse construction of the $23$ shifted residuals and their
sum of squares, including degree four, the unit normalization residual,
$346$ coefficient records and the twenty-term index bound.  A second
independent implementation agrees on all coefficient and bound records.
\item Direct T4 enumeration and independent class-factorized counting,
agreeing on all $37\,075$ ordered tiles, plus the support, carry and
population checks.  The direct enumeration examines $20\,624\,220$
quadruples after the horizontal head filter.
\item Exact polynomial identities for the four-factor nine-variable
norm and its unconditional majorant, small complete-homogeneous coding
identities, the factorial-weighted coding identity, and $432$ constructed
relation-combiner witnesses.
\item Sixty-four modular comparisons of generated nine-variable circuits
with separate direct formulas on smaller coding dimensions, followed by
$24$ comparisons using the actual fixed T4 and index coefficients.
The latter also compare $840$ named intermediate values.  Geometric
sums are evaluated without division, including at noninvertible radices.
\item A symbolic leading-coefficient pass through both nine-variable
circuits, with all $348$ named fixed leaves treated as independent
symbols.  The final leading coefficient is a nonzero scalar, so the
degree claim survives the authentic coefficient specialization.
\item Forty-eight comparisons of the twelve-variable circuit with an
independent eight-dimensional radical algebra, structural gate checks,
and two modular line-degree diagnostics.  Its exact degree additionally
has the uniform Laurent-series proof in the mathematical section.
\item For each signed weight, $96$ direct-formula comparisons, together
with $144$ isolated inverse-component checks, $6786$ finite gap-obstruction
checks and circuit/degree audits.  The signed proof establishes the
inequalities universally; these finite checks are supporting evidence.
\end{itemize}

For independent arithmetic fingerprints, the actual coefficient residues
and index residues are
\begin{center}
\small
\begin{tabular}{@{}rrrrrrr@{}}
\toprule
Modulus & $\beta$ & $K_0$ & $M_C$ & $M_F$ & $z_\star$ & $y_\star$\\\midrule
101 & 9 & 12 & 28 & 32 & 49 & 12\\
1009 & 1 & 1000 & 520 & 356 & 409 & 696\\
1000003 & 160640 & 704240 & 343080 & 140964 & 259269 & 170996\\
\bottomrule
\end{tabular}
\end{center}
For these prime moduli, $s_\star$ is first evaluated modulo $p-1$
using the actual ports modulo $p-1$, and Fermat's theorem then gives
$z_\star\bmod p$.  This is valid because $p$ is odd and $2$ is a unit.
No Euler/Fermat exponent reduction is silently used for a nonunit.
The bounded digit-polynomial evaluation then gives $y_\star\bmod p$.
These are checks of the literal constant definition, not a search for
a consistency witness or evidence of nonsolvability.

\subsection{Archive organization and replay}

The root \texttt{README.md} gives the command sequence and metric
conventions.  \texttt{report.tex}, the three included section files and
\texttt{report.pdf} form the complete paper.  The directories have these
roles:
\begin{center}
\begin{tabular}{@{}lp{10.2cm}@{}}
\toprule
Directory & Contents\\\midrule
\texttt{base/} & Exact source quartic; independently generated normalized coefficients and index bound.\\
\texttt{coefficients/} & Ordered T4 tiles and the frozen compiler coefficient recipe.\\
\texttt{code/} & Independent base constructor, actual fixed-port evaluator and replay support.\\
\texttt{nine/} & Two literal arithmetic DAGs, constructor, identity/degree checks and actual-instance modular verifier.\\
\texttt{direct12/} & Twelve-variable circuit, generator, independent checks and full proof.\\
\texttt{integer27/} & Both signed weights, product/system circuits, exact ledgers and full proof.\\
\texttt{proofs/} & Source-interface audit, independent review and comparison notes.\\
\texttt{upstream/} & Selected pinned mathematical dependencies, T4 checkers, machine data, notices and provenance.\\
\texttt{verification/} & Current-session receipts and the integrated replay record.\\
\bottomrule
\end{tabular}
\end{center}

From the extracted package, run \texttt{python3 reproduce.py} for the
arithmetic and binding checks; \texttt{python3 reproduce.py --full-t4}
also re-enumerates the local language.  Run \texttt{make pdf} for two
pdfLaTeX passes.  No network access, giant coefficient expansion or
proof-assistant installation is needed for these finite checks.
The SHA-256 manifest records checksums for the delivered files; it is not a
mathematical proof certificate.  Selected upstream source references
retain their historical commit metadata, while the acquisition and
comparison pin of this package is the full hash on the title page.

\subsection{Logical limits and next proof work}

No equation is asserted to have or lack a solution independently of
the requested equivalence.  Under consistency, G\"odel's second
incompleteness theorem prevents ZFC from proving the no-solution claim.
An assertion of two-sided independence would require additional care:
mere consistency does not generally rule out a theory proving its own
inconsistency.  No stronger independence conclusion is used here.

The most useful next formalization tasks are concrete: bind the
$32$-variable sparse polynomial and its coefficient bound to
\texttt{Jones1982.Nine.wset\_iff\_nine}; formalize the unconditional
majorant and two-radical finalizer; encode the twelve-variable uniqueness
and integrality lemmas; and formalize the signed denominator exclusions.
The finite coefficient adapter also needs a proof-assistant interface
before the entire new Diophantine endpoint can be called kernel checked.

For further simplification, the degree/circuit cost of the two separating
weights in the twelve-variable proof is a plausible target.  Another
target is a smaller normalized quartic, because the Jones coding exponent
grows as $5^{\nu+1}$.  Neither observation is a claimed result.  The
supplied equations, coefficient definitions and proofs are the concrete
deliverable.

\begin{thebibliography}{99}
\bibitem{jones1982}
James P. Jones, \emph{Universal Diophantine equation},
Journal of Symbolic Logic \textbf{47}(3) (1982), 549--571.
\href{https://doi.org/10.2307/2273588}{doi:10.2307/2273588}.

\bibitem{collection}
ProveIt contributors, \emph{Con(ZFC) in Diophantine form}, wave-2
collection and comparison, pinned repository source.
\repofile{SetTheory/Cardinals/docs/reports/hilbert-tenth-problem/con-zfc-diophantine-forms/README.md}{Collection README}.
Package numbers 09, 11, 12 and 14 throughout this paper refer to the
directories in this collection at \texttt{f0ae4e0991ca}.

\bibitem{k215ordinary}
ProveIt contributors, \emph{Ordinary proof of the combined215 machine}.
\repofile{Computability/BusyBeaver/Research/ZFC/candidate-k215/proof.md}{Candidate-K215 ordinary proof}.
Machine and inherited mathematical attributions are retained in the
bundled source notices.

\bibitem{k215formal}
ProveIt contributors, \emph{Formal correctness of the exact 215-state machine},
\repofile{Computability/BusyBeaver/Research/ZFC/k225-formal/verification/k215/README.md}{verification README};
\repofile{Computability/BusyBeaver/Lean/BusyBeaver/ZFC/KEndpoint215.lean}{Lean endpoint source}.
Published build claims are cited as source records and were not rerun here.

\bibitem{t4}
ProveIt contributors, package 11, \emph{Audit and smaller coefficient backend
for the combined consistency equation}.
\repofile{SetTheory/Cardinals/docs/reports/hilbert-tenth-problem/con-zfc-diophantine-forms/11-t4/bridge/BRIDGE_AUDIT.md}{T4 adapter proof};
\repofile{SetTheory/Cardinals/docs/reports/hilbert-tenth-problem/con-zfc-diophantine-forms/11-t4/bridge/check_t4.py}{exact local predicate and finite check}.

\bibitem{pellinterface}
ProveIt contributors, \emph{Fixed blank compiler proof}, \emph{Pell lemma audit},
and the imported theorem ledger retained in package 11.
\repofile{SetTheory/Cardinals/docs/reports/hilbert-tenth-problem/con-zfc-diophantine-forms/12-gap-norm14/upstream/fixed_blank_67_proof.md}{Fixed blank compiler proof};
\repofile{SetTheory/Cardinals/docs/reports/hilbert-tenth-problem/con-zfc-diophantine-forms/11-t4/bridge/sources/proof_dependencies.json}{dependency ledger}.

\bibitem{quartic32}
ProveIt contributors, package 09, \emph{Inputless ordinary-strong systems
and a 32-variable quartic}.
\repofile{SetTheory/Cardinals/docs/reports/hilbert-tenth-problem/con-zfc-diophantine-forms/09-combined13/proofs/LOW_DEGREE_PROOF.md}{Low-degree proof};
\repofile{SetTheory/Cardinals/docs/reports/hilbert-tenth-problem/con-zfc-diophantine-forms/09-combined13/low_degree/quartic32.json}{exact sparse quadratic system}.

\bibitem{root13}
ProveIt contributors, package 09, \emph{Root amalgamation proof} and
independent arithmetic review.
\repofile{SetTheory/Cardinals/docs/reports/hilbert-tenth-problem/con-zfc-diophantine-forms/09-combined13/proofs/ROOT_AMALGAMATION_PROOF.md}{Three-root amalgamation source}.

\bibitem{integer29}
ProveIt contributors, package 14, \emph{Signed-domain proof} and
unrestricted-integer constructions.
\repofile{SetTheory/Cardinals/docs/reports/hilbert-tenth-problem/con-zfc-diophantine-forms/14-integer29/proofs/signed_domain_proof.md}{Signed-core proof};
\repofile{SetTheory/Cardinals/docs/reports/hilbert-tenth-problem/con-zfc-diophantine-forms/14-integer29/README.md}{integer29 comparison}.

\bibitem{nineformal}
ProveIt contributors, \emph{Jones 1982, reduction to nine unknowns},
Lean sources in \path{Computability/HilbertTenthProblem/Lean/Diophantine}.
\repofile{Computability/HilbertTenthProblem/Lean/Diophantine/Paper1982/NineMain.lean}{NineMain};
\repofile{Computability/HilbertTenthProblem/Lean/Diophantine/Paper1982/NineReduce.lean}{NineReduce};
\repofile{Computability/HilbertTenthProblem/Lean/Diophantine/Paper1982/NineRatio.lean}{NineRatio};
\repofile{Computability/HilbertTenthProblem/Lean/Diophantine/Paper1982/NinePoly.lean}{NinePoly};
\repofile{Computability/HilbertTenthProblem/Lean/Diophantine/Paper1982/Coding.lean}{Coding};
\repofile{Computability/HilbertTenthProblem/Lean/Diophantine/Paper1982/Master1.lean}{Master1};
\repofile{Computability/HilbertTenthProblem/Lean/Diophantine/Paper1982/Index.lean}{Index}.
\end{thebibliography}
\end{document}
